\scanpage{84}

fixant les arguments $t_j$ pour $j\ne i$, $L_k(t)=L_k(t_1,\ldots t_n)$ \deleted{serait une fonction linéaire strictement décroissante de \deleted{\unsure{$t_i$}} $s=t_i$, majorée par la fonction continue croissante $\inserted{\varphi(t)=}\varphi(t_1,\ldots,t_n)$, \deleted{\unsure{donc}} ce qui est manifestement impossible, \deleted{\unsure{puisque}} car \unsure{$L$|$\lambda$}$(s)\to\infty$ si $s\to-\infty$, et que $\varphi(s)$ reste borné pour $s\to-\infty$.}

on aurait $L(t)\to\unsure{0|\infty}$ si $t_i\to-\infty$, tandis que $\varphi(t)$ \deleted{\unsure{restant|resterait}}\inserted{resterait} borné pour $t_i\to-\infty$, ce \deleted{\unsure{(en effet)}} qui est incompatible avec $L(t)\leq\varphi(t)$ pour \deleted{tout} $t$. \deleted{\unsure{$f$}} Comme \deleted{\unsure{maintenant}}\inserted{enfin} $\varphi$ est aussi supposée symétrique, $L\deleted{\unsure{(t_1,\ldots,t_n)}}\leq\varphi\deleted{\unsure{(t_1,\ldots,t_n)}}$ \inserted{($L$ fonction sur $R^n$)} implique $\deleted{\unsure{L_k(t_{\sigma1},\ldots,t_{\sigma n})\leq\varphi(t_{\sigma1},\ldots,t_{\sigma n})}}\inserted{\sigma L\leq\varphi}$ pour \deleted{tout} $\sigma\in\unsure{\mathscr S_n|\mathfrak S_n}$ (groupe des permutations des $n$ premiers entiers, i.e. \inserted{avec $\sigma L(t_1,\ldots t_n)=L(t_{\sigma^{-1}1},\ldots t_{\sigma^{-1}n})$}), donc on aura aussi

\[
\deleted{\unsure{\varphi(t_1,\ldots,t_n)|\varphi(t)}}\varphi=\sup_k\sup_{\sigma\in\mathscr S_n}\sigma.L_{k}\deleted{\unsure{(t_1,\ldots,t_n)}}\tag*{(7)}
\]

Il s'ensuit que \deleted{\unsure{si}} pour prouver (\unsure{1|3}), il suffit de prouver que \inserted{pour $L=L_k$, on a} \deleted{$\sigma L_k\deleted{\unsure{(\alpha_{\sigma1})}}(\alpha_1,\ldots\alpha_n)\leq\varphi(\beta_1,\ldots\beta_n)$ pour tout $k$, et tout $\sigma\in\mathscr S_n$. \unsure{Remplaçant} \deleted{\unsure{Désignant}} \unsure{enfin|alors} $\sigma L_k$ par $L$, on est ramené à prouver}

\[
\sup_{\sigma\in\mathscr S_n}(\sigma.L)(\alpha_1,\ldots\alpha_n)\leq\sup_{\sigma\in\mathscr S_n}(\unsure{\sigma|\tau}L)(\beta_{\unsure{\sigma1}},\ldots\beta_{\unsure{\sigma n}})\tag*{(8)}
\]

(car le deuxième mb en effet est $\leq\varphi(\beta_1,\ldots\beta_n)$) \deleted{\unsure{Pour le|Dans}}\inserted{Dans cette} inégalité, \deleted{\unsure{$L$ est une forme linéaire}} $L(t_1,\ldots t_n)=\sum_i^n a_it_i+b$, \deleted{\unsure{affine et}} avec des \deleted{\unsure{$a_i$}} coefficients quelconques dans le cas b), et des $a_i>0$ pour tout $i$ dans le cas a) du lemme. \deleted{\unsure{Désignons} \unsure{Remarquons},} On peut évidemment \deleted{\unsure{retrancher des deux membres} \unsure{l'effet de} \illegible{2}} supposer $b=0$.

\deleted{D'après un lemme connu, Il suffit \deleted{$\sum a_i b_i$} $L(\alpha_1,\ldots,\alpha_n)\leq L(\unsure{\beta_1,\ldots})$ \unsure{Notons que|lorsque} $\sum a_i\alpha_i\leq\sum a_i\alpha_{\sigma i}$ quand $\sigma$ est une permutation telle que la suite des $a_{\sigma i}$ est \uline{décroissante}, ce qui}

D'après un lemme connu, \inserted{valable du fait que $\alpha_1\geq\alpha_2\geq\ldots\geq\alpha_n$,} la somme $\sum a_{\sigma i}\alpha_i\inserted{=\sum a_i\alpha_{\sigma^{-1}i}}$ est maximum, pour $\sigma\in\mathscr S_n$ variable,

% archive sheet number: 85
